\documentclass[11pt]{article}
\usepackage[letterpaper,margin=0.95in]{geometry}
\usepackage[T1]{fontenc}
\usepackage{lmodern}
\usepackage{amsmath,amssymb,amsthm}
\usepackage{booktabs,array,longtable}
\usepackage{microtype,needspace}
\usepackage{xcolor}
\usepackage{fancyhdr}
\usepackage{hyperref}
\hypersetup{colorlinks=true,linkcolor=blue!45!black,citecolor=blue!45!black,
 urlcolor=blue!45!black,pdftitle={Closed Forms for Platonic Helices},
 pdfauthor={Robert L. Read}}
\pagestyle{fancy}
\fancyhf{}
\fancyhead[L]{\small Closed forms for Platonic helices}
\fancyhead[R]{\small Author-review draft}
\fancyfoot[C]{\thepage}
\setlength{\headheight}{14pt}
\setlength{\parskip}{0.35em}
\newtheorem{theorem}{Theorem}[section]
\newtheorem{proposition}[theorem]{Proposition}
\newtheorem{corollary}[theorem]{Corollary}
\theoremstyle{definition}
\newtheorem{definition}[theorem]{Definition}
\newcommand{\nb}{\mathbf n_B}
\newcommand{\nc}{\mathbf n_C}
\newcommand{\ee}{\mathbf e}
\newcommand{\uu}{\mathbf u}
\newcommand{\hh}{\widehat{\mathbf h}}
\newcommand{\Rot}{\operatorname{Rot}}
\newcommand{\tr}{\operatorname{tr}}
\newcommand{\eff}{\tau_{\mathrm{eff}}}
\newcommand{\degc}{^{\circ}}
\title{Closed Forms for Platonic Helices\\[0.4em]
 \large Exact parameters and an audit of the 28-species catalog}
\author{Robert L. Read\\\small Public Invention}
\date{October 2, 2026\\\small Working draft for author review; not submitted}

\begin{document}
\maketitle
\begin{abstract}
Repeated identical face-to-face placements of a rigid solid generate a screw
motion, and the centroids of successive joining faces trace a segmented helix.
For joint centers lying on equal-radius face normals, we express the principal
subunit rotation, the angle between a segment and the oriented screw axis, the
axial advance, and the joint-centroid radius in terms of three intrinsic inputs:
the normal-alignment angle, the joining twist, and the distance between joint
centers. A quaternion composition supplies a common denominator for all four
parameters. We specialize the result to the five Platonic solids and give exact
expressions for every entry in the previously published 28-species catalog,
with edge length one. An audit distinguishes the rotation of a subunit from the
geometry of a straight joint orbit, identifies the constructor-dependent length
normalization of the earlier table, and derives exact replacements for three
icosahedral twists that only approximately match face vertices. The original
and corrected configurations are both tabulated. Symbolic identities and
independent matrix, circumcircle, and vertex-matching checks accompany the
derivations.
\end{abstract}

\section{Introduction}
Identical rigid subunits conjoined by the same rule form a periodic chain.
The resulting motion is a screw motion, so a chosen point of a subunit traces
a discrete helical orbit. This observation connects the geometry of protein
helices, polyhedral chains, and engineered assemblies
\cite{lord2002,dauter2018,read2022}. Exact descriptions of the tetrahelix and
modified tetrahedral chains already exist, including exact rotation, radius,
and pitch calculations \cite{sadler2019,onaka2024}. The question addressed here
is how to obtain a uniform description from the geometry of the joining faces.

Read's 2022 paper supplied algorithms and a numerical catalog of 28
representative Platonic helices, and proposed closed-form descriptions as
future work \cite{read2022}. This paper develops such a description for the
\emph{helix through the joint-face centroids}. A vertex orbit or the orbit of
the center of the solid can have a different radius under the same screw
motion. Identifying the point whose orbit is being measured is therefore
essential when comparing results.

The main result uses the intrinsic inputs
\[
 L=\|C-B\|,\qquad
 \alpha=\arccos(-\nb\cdot\nc),\qquad \tau,
\]
where $B,C$ are selected joint centers within one subunit, $\nb,\nc$ are their
outward unit face normals, and $\tau$ is a twist measured from a specified
normal-alignment reference. The angle $\alpha$ is generally \emph{not} the
angle between the two outward normals. Nor should $\alpha,\tau$ be identified
without a derivation with the bond angle and torsion angle of a polygonal
chain used in other helix descriptions \cite{dauter2018}.

The theorem applies when the centers, relative to a common origin, satisfy
$B=\rho\nb$ and $C=\rho\nc$ for the same $\rho>0$. Every Platonic solid has
this property when its face centroids are used. Under this hypothesis,
\[
 Q=1-\cos^2(\alpha/2)\cos^2(\tau/2)
\]
governs all four requested parameters. The rotation formula has a broader
scope than the radius and advance formulas: face positions that do not
satisfy the equal-radius hypothesis require additional positional data.

Two distinctions emerge when the closed forms are compared with the earlier
catalog. First, a straight sequence of joint centers does not reveal the
rotation of the solid about that line. Second, the published lengths agree
with the original graphics constructors, which use unit circumradius for
four solids and unit edge length for the cube. We give an explicit conversion
to unit edges. A separate vertex audit finds that the three catalog twists
$45\degc,-75\degc,165\degc$ for one icosahedral face class are approximate
matching choices. Their exact parameters remain well defined, and we retain
them in the literal catalog tables. We also derive the exact matching twists
and their parameters.

The word ``species'' below refers to the representatives and names inherited
from the 2022 catalog. We do not infer uniqueness of an ordered discrete helix
from its radius and axial advance alone: its angular step also matters. We
also make no claim that every chain in the catalog avoids self-intersection.

\section{Geometry, motion, and conventions}\label{sec:setup}
\subsection{A periodic joint-centroid orbit}
Place a reference solid with its center at the origin and choose two distinct
faces. Let their centroids be $B,C$ and their outward unit normals be
$\nb,\nc$. For a Platonic solid of inradius $\rho$,
\begin{equation}\label{eq:centered}
 B=\rho\nb,\qquad C=\rho\nc.
\end{equation}
Here $L=\|C-B\|$ is a straight-line distance within one solid. It is also the
length of one segment of the joint-centroid helix. It is not the edge length,
the circumradius, or the advance along the screw axis.

Let $R$ be the forward rotation that makes the incoming $B$ face of the next
solid flush against the current $C$ face. Thus $R\nb=-\nc$. The rigid
placement is
\begin{equation}\label{eq:placement}
 T(\mathbf x)=R(\mathbf x-B)+C
            =R\mathbf x+\mathbf t,\qquad \mathbf t=C-RB.
\end{equation}
In particular $T(B)=C$. Repeating this placement gives
$\mathbf x_j=T^j(B)$, so $\mathbf x_0=B$, $\mathbf x_1=C$, and every
consecutive segment has length $L$.

For a nonidentity rotation, let $\hh$ be a directed unit vector along its
screw axis and let $\theta\in[0,\pi]$ be its principal rotation angle. We
define
\begin{equation}\label{eq:dphi}
 d=(C-B)\cdot\hh=\mathbf t\cdot\hh,
 \qquad \phi=\arccos\!\left(\frac dL\right)\in[0,\pi].
\end{equation}
The equality follows because $R\hh=\hh$, which implies
$\hh\cdot RB=\hh\cdot B$. The radius $r$ is the perpendicular distance
from $B$ (or any $\mathbf x_j$) to the screw-axis line.

If $r>0$, suitable perpendicular unit vectors $\mathbf e_1,\mathbf e_2$ and
a point $\mathbf a$ on the axis give
\begin{equation}\label{eq:orbit}
 \mathbf x_j=\mathbf a+
 r\{\cos(j\theta)\mathbf e_1+\sin(j\theta)\mathbf e_2\}
 +jd\hh,
 \qquad \mathbf e_2=\hh\times\mathbf e_1.
\end{equation}
This is a statement about an ordered joint orbit, rather than merely the
continuous cylinder or circle containing it.

\subsection{The zero-twist reference}\label{sec:reference}
For $0<\alpha<\pi$, define
\begin{equation}\label{eq:alphaeu}
 \alpha=\arccos(-\nb\cdot\nc),\qquad
 \ee=\frac{C-B}{L},\qquad
 \uu=\frac{\nb\times(-\nc)}{\sin\alpha}.
\end{equation}
Let $A$ be the shortest rotation that sends $\nb$ to $-\nc$. It has axis
$\uu$ and angle $\alpha$. Our zero-twist placement uses this rotation.
The additional twist is about the outward normal of the receiving face:
\begin{equation}\label{eq:composition}
 R=\Rot(\nc,\tau)A.
\end{equation}
For directed-axis formulas, first choose a representative
$-\pi<\tau<\pi$. The radius and principal rotation remain periodic in
$\tau$; the directed-axis choice can reverse when the representative crosses
the half-turn boundary.

This reference does not assert that $\tau=0$ matches the vertices of two
regular faces. Vertex matching depends on directions within the face plane
and is derived separately in Section~\ref{sec:matching}.

\begin{proposition}\label{prop:geometry}
Under~\eqref{eq:centered}, for distinct nonopposite faces,
\begin{equation}\label{eq:Lrho}
 L=2\rho\cos(\alpha/2),
 \qquad
 \ee=\cos(\alpha/2)\nc+
       \sin(\alpha/2)(\nc\times\uu),
 \qquad \ee\cdot\uu=0.
\end{equation}
\end{proposition}
\begin{proof}
Since $\nb\cdot\nc=-\cos\alpha$,
$\|\nc-\nb\|^2=2+2\cos\alpha=4\cos^2(\alpha/2)$.
This proves the first identity and gives
$\ee=(\nc-\nb)/(2\cos(\alpha/2))$. The vector triple-product identity gives
\[
 \nc\times\uu=
 \frac{-\nb-\cos\alpha\,\nc}{\sin\alpha}.
\]
Substitution and the half-angle identities prove the second identity.
The final identity follows because $\uu$ is perpendicular to both normals.
\end{proof}

\section{A common closed form for the four parameters}\label{sec:theorem}
\begin{theorem}\label{thm:main}
Suppose~\eqref{eq:centered} holds, $0<\alpha<\pi$, and
$-\pi<\tau<\pi$. Use the placement~\eqref{eq:placement} and reference
rotation~\eqref{eq:composition}. Define
\begin{equation}\label{eq:Q}
 Q(\alpha,\tau)=1-\cos^2(\alpha/2)\cos^2(\tau/2).
\end{equation}
Then $Q>0$, and the directed screw axis may be chosen as
\begin{equation}\label{eq:axis}
 \hh=\frac{\sin(\alpha/2)\cos(\tau/2)\uu+
                  \sin(\tau/2)\ee}{\sqrt Q}.
\end{equation}
For this choice, the joint-centroid helix has
\begin{equation}\label{eq:main}
 \boxed{\begin{aligned}
 \cos\theta&=1-2Q,
 &\theta&=\arccos(1-2Q),\\[2pt]
 \cos\phi&=\frac{\sin(\tau/2)}{\sqrt Q},
 &\phi&=\arccos\!\left(\frac{\sin(\tau/2)}{\sqrt Q}\right),\\[2pt]
 d&=\frac{L\sin(\tau/2)}{\sqrt Q},
 &r&=\frac{L\sin(\alpha/2)|\cos(\tau/2)|}{2Q}.
 \end{aligned}}
\end{equation}
The axis orientation agrees with the nonstraight rows of the earlier catalog:
negative representative twists give $d<0$ and $\phi>\pi/2$.
\end{theorem}

\begin{proof}
Write $a=\alpha/2$ and $t=\tau/2$ temporarily. A unit quaternion for $A$ is
$(\cos a,\sin a\,\uu)$ and one for the added twist is
$(\cos t,\sin t\,\nc)$. Their product has scalar part
$\cos a\cos t$ because $\nc\cdot\uu=0$, and vector part
\begin{align}
 \mathbf v&=\sin a\cos t\,\uu+
       \sin t\{\cos a\,\nc+\sin a(\nc\times\uu)\}\notag\\
           &=\sin a\cos t\,\uu+\sin t\,\ee.\label{eq:qvector}
\end{align}
The last equality is Proposition~\ref{prop:geometry}. Orthogonality gives
\[
 \|\mathbf v\|^2=\sin^2a\cos^2t+\sin^2t
                 =1-\cos^2a\cos^2t=Q.
\]
Since the scalar part is nonnegative on the chosen twist interval, this is
the principal-angle quaternion. Therefore
\begin{equation}\label{eq:halfangle}
 \sin(\theta/2)=\sqrt Q,
 \qquad \cos(\theta/2)=\cos(\alpha/2)\cos(\tau/2),
\end{equation}
and its vector part divided by $\sqrt Q$ is the directed axis~\eqref{eq:axis}.
This proves the rotation formula.

Taking the scalar product of~\eqref{eq:axis} with $C-B=L\ee$ gives
$d=L\sin(\tau/2)/\sqrt Q$. Dividing by $L$ proves the formula for
$\cos\phi$. No choice of the position of an axis point is needed for this
calculation: the axial component of the displacement is the same at every
point of a rigid screw motion.

Project $B,C$ orthogonally into a plane perpendicular to the axis. Their
projected chord length $c$ satisfies
\begin{align}
 c^2&=L^2-d^2
     =L^2\frac{Q-\sin^2t}{Q}
     =L^2\frac{\sin^2a\cos^2t}{Q}.
\end{align}
The chord of a circle of radius $r$ subtending $\theta$ is
\begin{equation}\label{eq:chord}
 c=2r\sin(\theta/2).
\end{equation}
Combining this with~\eqref{eq:halfangle} proves the radius formula.
\end{proof}

The four outputs are linked by the familiar identities
\begin{equation}\label{eq:linked}
 d=L\cos\phi,\qquad
 r=\frac{L\sin\phi}{2\sin(\theta/2)}.
\end{equation}
The geometric content of the theorem is the explicit connection from face
data $(\alpha,\tau,L)$ to these extrinsic helix parameters.

\begin{corollary}\label{cor:shape}
The quantities $\theta,\phi,r/L,d/L$ depend on $\alpha,\tau$ alone. Scaling
the solid by a factor $s>0$ scales $L,r,d,\rho$ by $s$ and leaves both angles
unchanged. For nonzero twist,
\begin{equation}\label{eq:tanphi}
 \tan^2\phi=\sin^2(\alpha/2)\cot^2(\tau/2).
\end{equation}
\end{corollary}
\begin{proof}
The scale claim follows from~\eqref{eq:main}. The final identity is
$(L^2-d^2)/d^2$, using the two squared expressions in the proof.
The squared identity does not itself specify the sign of $\cos\phi$.
\end{proof}

\subsection{A different twist reference}
Let $R_0$ be any other zero-twist rotation satisfying $R_0\nb=-\nc$.
Then $R_0A^{-1}$ fixes $\nc$, so for some phase $\gamma$,
\begin{equation}\label{eq:phase}
 R_0=\Rot(\nc,\gamma)A,
 \qquad
 \eff=\tau+\gamma\pmod{2\pi}.
\end{equation}
The principal rotation law is always
\begin{equation}\label{eq:generalrotation}
 \cos(\theta/2)=\cos(\alpha/2)|\cos(\eff/2)|.
\end{equation}
When~\eqref{eq:centered} holds, replacing $\tau$ by a principal representative
of $\eff$ in Theorem~\ref{thm:main} also supplies $d,\phi,r$. Without the
centroid-position hypothesis, the rotation law still holds but the remaining
three formulas need not hold.

This is relevant to the repository convention. Its forward baseline is the
product of two shortest rotations
\begin{equation}\label{eq:sourcebaseline}
 R_0=S(\ee,\nc)S(\nb,-\ee),
\end{equation}
where $S(\mathbf a,\mathbf b)$ denotes a shortest alignment. Under
\eqref{eq:centered}, both rotations have axis $\uu$ and angle $\alpha/2$,
because
\[
 \nb\cdot(-\ee)=\ee\cdot\nc=\cos(\alpha/2).
\]
Their product is $A$. Thus $\gamma=0$ for the Platonic face-centroid cases,
and the original tabulated twists can be used directly. This conclusion is
derived from the geometry, rather than selected to force numerical agreement.

\section{Exact face geometry for unit edges}\label{sec:faces}
Throughout the remaining tables the edge length is one. For another edge
length $a$, multiply every tabulated $L,r,d$ by $a$. Define the golden ratio
using a symbol distinct from the helix angle:
\begin{equation}\label{eq:constants}
 w=\sqrt5,\qquad g=\frac{1+w}{2},\qquad
 K=\frac{\sqrt{5+2w}}{2w},\qquad
 p=\sqrt6+\sqrt2,\qquad q=\sqrt6-\sqrt2.
\end{equation}
These are fixed algebraic constants. In particular, $K$ is the unit-edge
Dodecagear radius; it is not an additional input to the theorem.

\begin{center}
\renewcommand{\arraystretch}{1.35}
\begin{tabular}{@{}lcc@{}}
\toprule
Solid & Inradius $\rho$ & Circumradius for edge one\\
\midrule
Tetrahedron & $\sqrt6/12$ & $\sqrt6/4$\\
Cube & $1/2$ & $\sqrt3/2$\\
Octahedron & $\sqrt6/6$ & $1/\sqrt2$\\
Dodecahedron & $\sqrt{250+110w}/20$ & $\sqrt3\,g/2$\\
Icosahedron & $g^2/(2\sqrt3)$ & $\sqrt{10+2w}/4$\\
\bottomrule
\end{tabular}
\end{center}

Let $A$ denote edge-adjacent faces, $S,T,F$ successive nonopposite normal
classes where needed, and $O$ opposite faces. The labels $S,T,F$ are defined
by the scalar products in the next table, rather than by an assumed graphical
ordering. Counts are relative to a fixed incoming face and exclude that face.

\begin{center}
\small\renewcommand{\arraystretch}{1.4}
\begin{tabular}{@{}llrcc@{}}
\toprule
Solid & Class & Faces & $\cos\alpha=-\nb\cdot\nc$ & $L$\\
\midrule
Tet & $A$ & 3 & $1/3$ & $1/3$\\
Cub & $A$ & 4 & $0$ & $1/\sqrt2$\\
    & $O$ & 1 & $1$ & $1$\\
Oct & $A$ & 3 & $-1/3$ & $\sqrt2/3$\\
    & $S$ & 3 & $1/3$ & $2/3$\\
    & $O$ & 1 & $1$ & $\sqrt6/3$\\
Dod & $A$ & 5 & $-1/w$ & $g^2/w$\\
    & $S$ & 5 & $1/w$ & $g^3/w$\\
    & $O$ & 1 & $1$ & $\sqrt{250+110w}/10$\\
Ico & $A$ & 3 & $-w/3$ & $g/3$\\
    & $S$ & 6 & $-1/3$ & $g^2/3$\\
    & $T$ & 6 & $1/3$ & $\sqrt2\,g^2/3$\\
    & $F$ & 3 & $w/3$ & $g^3/3$\\
    & $O$ & 1 & $1$ & $g^2/\sqrt3$\\
\bottomrule
\end{tabular}
\end{center}

\subsection{Coordinates and scalar products}
Unit-edge tetrahedral vertices can be taken as
$(\pm1,\pm1,\pm1)/(2\sqrt2)$ with product of the three signs positive.
The normals of distinct faces have dot product $-1/3$. Alternatively, if
two faces omit vertices $v_i,v_j$, subtracting their centroids gives
$(v_j-v_i)/3$, proving $L=1/3$ directly.

The cube has vertices $(\pm1/2,\pm1/2,\pm1/2)$ and normals along the
coordinate axes. The octahedron has vertices on those axes at distance
$1/\sqrt2$ from the origin. Its face centers are
$(\pm1,\pm1,\pm1)/(3\sqrt2)$. Changing $k$ signs from a reference face gives
\begin{equation}\label{eq:octgeometry}
 L_k=\frac{\sqrt{2k}}3,\qquad \cos\alpha_k=\frac{2k-3}{3},
 \qquad k=1,2,3.
\end{equation}

Dodecahedral unit-edge vertices can be taken as $g/2$ times
\[
 (\pm1,\pm1,\pm1),\quad
 (0,\pm g^{-1},\pm g),\quad
 (\pm g^{-1},\pm g,0),\quad
 (\pm g,0,\pm g^{-1}).
\]
The face normals point to the vertices of the dual icosahedron. Dividing
\[
 (0,\pm1,\pm g),\quad(\pm1,\pm g,0),\quad(\pm g,0,\pm1)
\]
by $\sqrt{1+g^2}$ gives their unit directions. Relative to one of these
directions, the other eleven have dot products $1/w,-1/w,-1$ with counts
$5,5,1$. For example,
$g/(1+g^2)=1/w$. The table then follows from~\eqref{eq:Lrho}.

Icosahedral unit-edge vertices are half the twelve coordinate vectors just
listed. Their face normals are the directions of a dual dodecahedron:
eight vectors $(\pm1,\pm1,\pm1)/\sqrt3$ and twelve vectors formed from
$(0,\pm g,\pm g^{-1})/\sqrt3$ by cyclic coordinate permutations. Relative
to $(-1,1,1)/\sqrt3$, the other nineteen have dot products
\[
 w/3,\quad 1/3,\quad -1/3,\quad -w/3,\quad -1
\]
with counts $3,6,6,3,1$. Negating these dot products gives the listed
$\cos\alpha$. Computing distances from the origin to the face planes of
these coordinate solids gives the inradii above; thus the length table uses
the same edge normalization as the vertex coordinates.

\section{Circular, straight, and singular joint orbits}\label{sec:degenerate}
\subsection{Zero twist and the circular limit}
For nonopposite faces, putting $\tau=0$ gives
\begin{equation}\label{eq:torus}
 d=0,\qquad \phi=\pi/2,\qquad \theta=\alpha,
 \qquad r=\frac{L}{2\sin(\alpha/2)}.
\end{equation}
The joint orbit lies on a circle. It closes after finitely many steps exactly
when $\theta/(2\pi)$ is rational. A catalog name ending in ``torus'' does
not itself establish finite closure of the ordered orbit or of the solid chain.

As $\tau$ runs from $0$ to $\pi$, the formulas give $d$ increasing from zero
to $L$ and $r$ decreasing to zero. For example, writing $s=\sin(\tau/2)$,
\[
 \frac dL=\frac{s}{\sqrt{\sin^2(\alpha/2)+\cos^2(\alpha/2)s^2}},
\]
whose derivative with respect to $s$ is positive. Similarly, the expression
for $r$ is increasing as a function of $|\cos(\tau/2)|$. These statements
describe the continuous family of joining twists, not only the matching
twists selected for a catalog.

\subsection{A straight orbit can hide a rotation}
At $\tau=\pm\pi$ and nonopposite faces, the continuous formulas give
$r=0$, $|d|=L$, and $\theta=\pi$. The points lie on the screw axis even
though the solid undergoes a half-turn. The limiting directed axes for
$+\pi$ and $-\pi$ are opposite. In all straight catalog rows we choose
the joint-line direction from $B$ to $C$, record $d=L$, and set $\phi=0$.

For opposite faces $\alpha=0$, $A=I$ and $\ee=\nc$. If the twist is
nonzero, the screw axis is the line through both face centers. Hence $r=0$
and the same forward convention gives $d=L,\phi=0$. Direct substitution
into the directed-axis expression would instead give
$\cos\phi=\operatorname{sgn}(\sin(\tau/2))$ for a principal nonzero twist.
The difference is an axis-direction choice, not a different line of joint
points. The full subunit rotation is still the principal angle of the twist.

The earlier catalog records $\theta=0$ for each straight point sequence.
Here $\theta$ always denotes the principal rotation of the specified rigid
placement. The two conventions differ in precisely these five rows:
\begin{center}
\renewcommand{\arraystretch}{1.3}
\begin{tabular}{@{}lrrr@{}}
\toprule
Name & Twist (deg.) & Subunit $\theta$ (deg.) & Catalog marker (deg.)\\
\midrule
Boxbeam & $-180$ & $180$ & $0$\\
Staircase & $-180$ & $180$ & $0$\\
Octabeam & $-60$ & $60$ & $0$\\
Dodecabeam & $-108$ & $108$ & $0$\\
Pearlshaft & $-60$ & $60$ & $0$\\
\bottomrule
\end{tabular}
\end{center}
The angular step is unidentifiable from collinear joint points alone. This
does not imply that the specified motion of a labeled subunit has no rotation.

\subsection{Pure translation and a repeated joint center}
For opposite faces and zero twist, $Q=0$ and the motion is pure translation.
The formulas with $Q$ in the denominator must not be used. The joint orbit
is a line; we record its natural joint-line representation as
$r=0,d=L,\phi=0,\theta=0$. A pure translation does not select a unique
axis line among parallel lines.

Selecting the same face twice gives $B=C$ and $L=0$, so the joint orbit
is stationary. The shortest normal alignment is then antipodal and its axis
is not uniquely specified by the normals. Such same-face constructions are
excluded here, as they are from the displayed 28-row catalog. A stationary
joint orbit by itself does not establish that the entire rigid motion is the
identity.

\section{Length normalization in the earlier catalog}\label{sec:scale}
The earlier table does not state a uniform length normalization in its
caption or column headings. Its entries agree with the geometries built by
the associated THREE.js constructors \cite{repository}. The tetrahedron,
octahedron, dodecahedron, and icosahedron constructors use radius one, meaning
circumradius one. The cube constructor instead gives edge length one.

If $a_0$ is the original constructor's edge length, the conversion is
\begin{equation}\label{eq:conversion}
 r_{\text{edge }1}=\frac{r_{\mathrm{catalog}}}{a_0},
 \qquad d_{\text{edge }1}=\frac{d_{\mathrm{catalog}}}{a_0}.
\end{equation}
Use underlying values before rounding when possible. The factors are
\begin{center}
\renewcommand{\arraystretch}{1.35}
\begin{tabular}{@{}lcc@{}}
\toprule
Solid & Constructor edge $a_0$ & Catalog-to-unit-edge multiplier\\
\midrule
Tet & $\sqrt{8/3}$ & $\sqrt6/4$\\
Cub & $1$ & $1$\\
Oct & $\sqrt2$ & $1/\sqrt2$\\
Dod & $2/(\sqrt3\,g)$ & $\sqrt3\,g/2$\\
Ico & $4/\sqrt{10+2w}$ & $\sqrt{10+2w}/4$\\
\bottomrule
\end{tabular}
\end{center}
The code first scales a constructor geometry by $s$ to obtain a requested
joint-center distance. Its factor named \texttt{SCALE\_TO\_UNIT\_EDGE} is
$1/s$, which restores the original constructor dimensions. To obtain unit
edges, the factor would instead be $1/(sa_0)$. Thus the normalization issue
is localized and can be corrected by scaling the length columns. It does not
change $\theta,\phi$ or the scale-independent formulas.

For instance, the tetrahedron with circumradius one has
$a_0=\sqrt{8/3}$ and $L=2\sqrt6/9$. At twist $-120\degc$ it has
\[
 r=\frac{\sqrt2}{15},\qquad d=-\frac2{\sqrt{15}},
\]
which round to $0.094,-0.516$ as in the catalog. At edge length one, the
same formulas instead give $r=\sqrt3/30$ and $d=-1/\sqrt{10}$.

Uniform edge normalization can create equal $(r,d)$ pairs across different
solids. For example, Tetratorus, Octagear, and Icosa Tree Star all have
$r=\sqrt3/6,d=0$ at edge length one, but different angular steps. Their
ordered joint orbits should therefore not be merged using only $(r,d)$.

\section{Exact parameters for the 28 original catalog inputs}\label{sec:catalog}
The tables in this section retain every original name, selected $C$-face,
and twist from Table~1 of \cite{read2022}. Face $0$ is the incoming $B$ face;
the numerical face labels follow the pinned repository implementation.
Lengths use edge length one, and $w,g,K,p,q$ are the fixed constants in
\eqref{eq:constants}. All square roots denote their nonnegative real values.
Angles are recovered by $\theta=\arccos(\cos\theta)$ and
$\phi=\arccos(\cos\phi)$. Numeric twists in these tables are degrees.

The dagger on $\cos\phi$ denotes the forward-axis convention for straight
joint rows. Their $\cos\theta$ values describe the full rigid placement,
not the earlier table's zero marker. The three class-$T$ icosahedral twists
remain literal catalog inputs in this section; exact vertex-matching
replacements appear in Section~\ref{sec:correction}. Thus a successful
parameter comparison is not being used as evidence of vertex matching.

\subsection{Tetrahedral and cubic entries}
For a tetrahedron $L=1/3$, $\cos\alpha=1/3$, and the radius specialization is
\begin{equation}\label{eq:tetra}
 r_{\mathrm{Tet}}(\tau)=
 \frac{\sqrt3 L|\cos(\tau/2)|}{2(1+2\sin^2(\tau/2))}.
\end{equation}
This expression is tetrahedron-specific; the face-angle constants change
for other solids. The table includes the two tetrahedral representatives.
\begin{center}
\small
\renewcommand{\arraystretch}{1.7}
\begin{tabular}{@{}lrrcccc@{}}
\toprule
Name & $C$-face & $\tau$ (deg.) & $\cos\theta$ & $\cos\phi$ & $r$ & $d$\\
\midrule
Tetrahelix & 1 & -120 & $\displaystyle - \frac{2}{3}$ & $\displaystyle - \frac{3}{\sqrt{10}}$ & $\displaystyle \frac{\sqrt{3}}{30}$ & $\displaystyle - \frac{1}{\sqrt{10}}$\\
Tetratorus & 1 & 0 & $\displaystyle \frac{1}{3}$ & $\displaystyle 0$ & $\displaystyle \frac{\sqrt{3}}{6}$ & $\displaystyle 0$\\
\bottomrule
\end{tabular}
\end{center}
For the cube, face $1$ is opposite face $0$ and face $2$ is adjacent.
Their lengths are $1$ and $1/\sqrt2$, respectively.
\begin{center}
\small
\renewcommand{\arraystretch}{1.7}
\begin{tabular}{@{}lrrcccc@{}}
\toprule
Name & $C$-face & $\tau$ (deg.) & $\cos\theta$ & $\cos\phi$ & $r$ & $d$\\
\midrule
Boxbeam & 1 & -180 & $\displaystyle -1$ & $\displaystyle 1^{\dagger}$ & $\displaystyle 0$ & $\displaystyle 1$\\
Staircase & 2 & -180 & $\displaystyle -1$ & $\displaystyle 1^{\dagger}$ & $\displaystyle 0$ & $\displaystyle \frac{1}{\sqrt{2}}$\\
Blockhelix & 2 & -90 & $\displaystyle - \frac{1}{2}$ & $\displaystyle - \sqrt{\frac{2}{3}}$ & $\displaystyle \frac{\sqrt{2}}{6}$ & $\displaystyle - \frac{1}{\sqrt{3}}$\\
Cubatorus & 2 & 0 & $\displaystyle 0$ & $\displaystyle 0$ & $\displaystyle \frac{1}{2}$ & $\displaystyle 0$\\
\bottomrule
\end{tabular}
\end{center}

\clearpage
\subsection{Octahedral entries}
For edge-sharing faces $L=\sqrt2/3$ and $\cos\alpha=-1/3$; for faces
sharing only a vertex $L=2/3$ and $\cos\alpha=1/3$. Consequently
\begin{equation}\label{eq:octfunctions}
 r_{\mathrm{Oct}}(\tau)=
 \begin{cases}
 \displaystyle\frac{\sqrt3|\cos(\tau/2)|}{3(2+\sin^2(\tau/2))},&A,\\[5pt]
 \displaystyle\frac{\sqrt3|\cos(\tau/2)|}{3(1+2\sin^2(\tau/2))},&S,\\[5pt]
 0,&O.
 \end{cases}
\end{equation}
Repository faces $1,3,7$ are class $A$, faces $2,4,6$ are class $S$, and
face $5$ is opposite. The nonopposite matching twists are
$0,\pm120\degc$; opposite matching twists are $-60,60,180\degc$.
\begin{center}
\small
\renewcommand{\arraystretch}{1.7}
\begin{tabular}{@{}lrrcccc@{}}
\toprule
Name & $C$-face & $\tau$ (deg.) & $\cos\theta$ & $\cos\phi$ & $r$ & $d$\\
\midrule
Octabeam & 5 & -60 & $\displaystyle \frac{1}{2}$ & $\displaystyle 1^{\dagger}$ & $\displaystyle 0$ & $\displaystyle \frac{\sqrt{6}}{3}$\\
Octaspiky & 1 & -120 & $\displaystyle - \frac{5}{6}$ & $\displaystyle - \frac{3}{\sqrt{11}}$ & $\displaystyle \frac{2 \sqrt{3}}{33}$ & $\displaystyle \frac{- \sqrt{22}}{11}$\\
Octamedium & 2 & -120 & $\displaystyle - \frac{2}{3}$ & $\displaystyle - \frac{3}{\sqrt{10}}$ & $\displaystyle \frac{\sqrt{3}}{15}$ & $\displaystyle \frac{- \sqrt{10}}{5}$\\
Octagear & 1 & 0 & $\displaystyle - \frac{1}{3}$ & $\displaystyle 0$ & $\displaystyle \frac{\sqrt{3}}{6}$ & $\displaystyle 0$\\
Treestar & 2 & 0 & $\displaystyle \frac{1}{3}$ & $\displaystyle 0$ & $\displaystyle \frac{\sqrt{3}}{3}$ & $\displaystyle 0$\\
\bottomrule
\end{tabular}
\end{center}

\clearpage
\subsection{Dodecahedral entries}
Faces $1$ and $2$ represent classes $A$ and $S$, and face $8$ is opposite.
The nonopposite matching twists are $0,\pm72\degc,\pm144\degc$.
The two continuous angle functions are particularly compact:
\begin{equation}\label{eq:dodphi}
 \begin{aligned}
 \cos\phi_A(\tau)&=
 \frac{\sqrt{10}\sin(\tau/2)}{\sqrt{10-(5-w)\cos^2(\tau/2)}},\\[4pt]
 \cos\phi_S(\tau)&=
 \frac{\sqrt{10}\sin(\tau/2)}{\sqrt{10-(5+w)\cos^2(\tau/2)}}.
 \end{aligned}
\end{equation}
The exact radius formulas below use $K$, the fixed radical in
\eqref{eq:constants}, to avoid long rationalized denominators.
\begin{center}
\small
\renewcommand{\arraystretch}{1.55}
\begin{tabular}{@{}lrrcc@{}}
\toprule
Name & $C$-face & $\tau$ (deg.) & $\cos\theta$ & $\cos\phi$\\
\midrule
Dodecabeam & 8 & -108 & $\displaystyle \frac{1 - w}{4}$ & $\displaystyle 1^{\dagger}$\\
Dodecadoubler & 1 & -144 & $\displaystyle \frac{- 2 w - 5}{10}$ & $\displaystyle - \sqrt{\frac{10}{13 - w}}$\\
The Alternater & 2 & -144 & $\displaystyle \frac{- w - 15}{20}$ & $\displaystyle - \sqrt{\frac{10}{17 - 3 w}}$\\
Dodecashaft & 1 & -72 & $\displaystyle \frac{w - 15}{20}$ & $\displaystyle - \sqrt{\frac{10}{3 w + 17}}$\\
Dodecagear & 1 & 0 & $\displaystyle \frac{- w}{5}$ & $\displaystyle 0$\\
Dodecacorkscrew & 2 & -72 & $\displaystyle \frac{2 w - 5}{10}$ & $\displaystyle - \sqrt{\frac{10}{w + 13}}$\\
Dodecadonut & 2 & 0 & $\displaystyle \frac{w}{5}$ & $\displaystyle 0$\\
\bottomrule
\end{tabular}
\end{center}
\begin{center}
\small
\renewcommand{\arraystretch}{1.6}
\begin{tabular}{@{}lcc@{}}
\toprule
Name & $r$ & $d$\\
\midrule
Dodecabeam & $\displaystyle 0$ & $\displaystyle \frac{\sqrt{110 w + 250}}{10}$\\
Dodecadoubler & $\displaystyle \frac{2 K w}{2 w + 15}$ & $\displaystyle -\frac{g^2}{w}\sqrt{\frac{10}{13 - w}}$\\
The Alternater & $\displaystyle \frac{4 K w}{w + 35}$ & $\displaystyle -\frac{g^3}{w}\sqrt{\frac{10}{17 - 3 w}}$\\
Dodecashaft & $\displaystyle \frac{4 K g^{2} w}{35 - w}$ & $\displaystyle -\frac{g^2}{w}\sqrt{\frac{10}{3 w + 17}}$\\
Dodecagear & $\displaystyle K$ & $\displaystyle 0$\\
Dodecacorkscrew & $\displaystyle \frac{2 K g^{2} w}{15 - 2 w}$ & $\displaystyle -\frac{g^3}{w}\sqrt{\frac{10}{w + 13}}$\\
Dodecadonut & $\displaystyle K g^{2}$ & $\displaystyle 0$\\
\bottomrule
\end{tabular}
\end{center}
For example, Dodecashaft uses $\tau=-72\degc$ and class $A$. Inserting
$\sin^2 36\degc=(5-w)/8$ and $\cos^2 36\degc=(3+w)/8$ gives
$\cos\phi=-\sqrt{10/(17+3w)}$. Multiplication by $L=g^2/w$ gives its
tabulated advance. Such factorization is used throughout the length tables.

\clearpage
\subsection{Icosahedral entries}\label{sec:icotables}
The representative faces $1,2,8,12,13$ have
\[
 \cos\alpha=-w/3,-1/3,1/3,w/3,1,
\]
respectively. In class $T$, $L=\sqrt2\,g^2/3$ and
\begin{equation}\label{eq:icophiT}
 \cos\phi_T(\tau)=
 \frac{\sqrt3\sin(\tau/2)}{\sqrt{1+2\sin^2(\tau/2)}}.
\end{equation}
The three original twists for this class are $165\degc,-75\degc,45\degc$.
Here $p,q$ abbreviate the radicals arising from the exact sine and cosine
of $15\degc$ and $75\degc$. The tables describe these precise input angles,
although they are not exact vertex matches.
\begin{center}
\small
\renewcommand{\arraystretch}{1.55}
\begin{tabular}{@{}lrrcc@{}}
\toprule
Name & $C$-face & $\tau$ (deg.) & $\cos\theta$ & $\cos\phi$\\
\midrule
Pearlshaft & 13 & -60 & $\displaystyle \frac{1}{2}$ & $\displaystyle 1^{\dagger}$\\
Quasi-planar & 8 & 165 & $\displaystyle \frac{- p - 2}{6}$ & $\displaystyle \sqrt{\frac{3 \left(p + 4\right)}{2 \left(p + 8\right)}}$\\
Two Strands & 1 & 120 & $\displaystyle \frac{- w - 9}{12}$ & $\displaystyle \sqrt{\frac{18}{w + 21}}$\\
Slow Twist & 12 & 120 & $\displaystyle \frac{w - 9}{12}$ & $\displaystyle \sqrt{\frac{18}{21 - w}}$\\
Rock Candy & 2 & 120 & $\displaystyle - \frac{5}{6}$ & $\displaystyle \frac{3}{\sqrt{11}}$\\
Icosa Tree Star & 1 & 0 & $\displaystyle \frac{- w}{3}$ & $\displaystyle 0$\\
Icosacorkscrew & 8 & -75 & $\displaystyle \frac{q - 2}{6}$ & $\displaystyle - \sqrt{\frac{3 \left(4 - q\right)}{2 \left(8 - q\right)}}$\\
Planar point cluster & 2 & 0 & $\displaystyle - \frac{1}{3}$ & $\displaystyle 0$\\
Big Icosacorkscrew & 8 & 45 & $\displaystyle \frac{-1 + \sqrt{2}}{3}$ & $\displaystyle \sqrt{\frac{9 - 3 \sqrt{2}}{14}}$\\
The Wheel & 12 & 0 & $\displaystyle \frac{w}{3}$ & $\displaystyle 0$\\
\bottomrule
\end{tabular}
\end{center}
\begin{center}
\small
\renewcommand{\arraystretch}{1.6}
\begin{tabular}{@{}lcc@{}}
\toprule
Name & $r$ & $d$\\
\midrule
Pearlshaft & $\displaystyle 0$ & $\displaystyle \frac{g^{2}}{\sqrt{3}}$\\
Quasi-planar & $\displaystyle \frac{g^{2} \sqrt{4 - p}}{\sqrt{3} \left(p + 8\right)}$ & $\displaystyle g^{2} \sqrt{\frac{p + 4}{3 \left(p + 8\right)}}$\\
Two Strands & $\displaystyle \frac{2 g^{2}}{\sqrt{3} \left(w + 21\right)}$ & $\displaystyle g \sqrt{\frac{2}{w + 21}}$\\
Slow Twist & $\displaystyle \frac{2 g^{2}}{\sqrt{3} \left(21 - w\right)}$ & $\displaystyle g^{3} \sqrt{\frac{2}{21 - w}}$\\
Rock Candy & $\displaystyle \frac{\sqrt{6} g^{2}}{33}$ & $\displaystyle \frac{g^{2}}{\sqrt{11}}$\\
Icosa Tree Star & $\displaystyle \frac{\sqrt{3}}{6}$ & $\displaystyle 0$\\
Icosacorkscrew & $\displaystyle \frac{g^{2} \sqrt{q + 4}}{\sqrt{3} \left(8 - q\right)}$ & $\displaystyle - g^{2} \sqrt{\frac{4 - q}{3 \left(8 - q\right)}}$\\
Planar point cluster & $\displaystyle \frac{g^{2}}{2 \sqrt{6}}$ & $\displaystyle 0$\\
Big Icosacorkscrew & $\displaystyle \frac{g^{2} \sqrt{17 \sqrt{2} + 26}}{14 \sqrt{3}}$ & $\displaystyle g^{2} \sqrt{\frac{3 - \sqrt{2}}{21}}$\\
The Wheel & $\displaystyle \frac{g^{4}}{2 \sqrt{3}}$ & $\displaystyle 0$\\
\bottomrule
\end{tabular}
\end{center}

\section{Analytical face matching and the icosahedral correction}\label{sec:matching}
\subsection{A face-plane calculation}
Suppose the two selected faces are congruent regular $m$-gons. Choose a
vertex $P_B$ of the incoming face and a vertex $P_C$ of the receiving face.
After the shortest normal alignment $A$, put
\begin{equation}\label{eq:outvectors}
 \mathbf a=A(P_B-B),\qquad \mathbf b=P_C-C.
\end{equation}
These have equal length and lie in the plane perpendicular to $\nc$.
The twist that aligns the selected vertices is
\begin{equation}\label{eq:betarule}
 \beta=\operatorname{atan2}\!\left(
 \nc\cdot(\mathbf a\times\mathbf b),\ \mathbf a\cdot\mathbf b\right).
\end{equation}
Matching this one vertex matches the entire regular polygon as an unordered
vertex set: all its vertices are equally spaced on the same circumcircle.
Every matching twist is therefore
\begin{equation}\label{eq:matchingrule}
 \tau=\beta+\frac{2\pi j}{m}\pmod{2\pi},\qquad j=0,\ldots,m-1.
\end{equation}
Choosing a different receiving vertex changes $\beta$ by a multiple of
$2\pi/m$ and hence gives the same set of matching twists.

Substitution of the coordinate solids in Section~\ref{sec:faces} gives the
following bases. The $C$-face numbering is the same as in the catalog.
\begin{center}
\small\renewcommand{\arraystretch}{1.35}
\begin{tabular}{@{}llcl@{}}
\toprule
Solid & Faces or class & Base twist & Increment\\
\midrule
Tet & all three other faces & $0$ & $120\degc$\\
Cub & all five other faces & $0$ & $90\degc$\\
Oct & nonopposite & $0$ & $120\degc$\\
    & opposite, face $5$ & $60\degc$ & $120\degc$\\
Dod & nonopposite & $0$ & $72\degc$\\
    & opposite, face $8$ & $36\degc$ & $72\degc$\\
Ico & $A,S,F$ & $0$ & $120\degc$\\
    & opposite, face $13$ & $60\degc$ & $120\degc$\\
    & $T$: faces $8,11,19$ & $\beta$ & $120\degc$\\
    & $T$: faces $9,10,17$ & $-\beta$ & $120\degc$\\
\bottomrule
\end{tabular}
\end{center}
The last two rows use the exact angle derived next. A base of $120\degc$
in the original implementation is equivalent to zero modulo the triangular
increment; a base of $45\degc$ is not equivalent to the derived $\beta$.

\subsection{An exact icosahedral base angle}\label{sec:correction}
Use the unit-edge icosahedral faces numbered $0$ and $8$. With the standard
coordinates used in the verifier, their vertices are
\begin{align*}
 F_B&=\{(-g/2,0,1/2),\ (0,1/2,g/2),\ (-1/2,g/2,0)\},\\
 F_C&=\{(0,1/2,-g/2),\ (-g/2,0,-1/2),\ (0,-1/2,-g/2)\}.
\end{align*}
Thus
\begin{equation}\label{eq:icoBC}
 B=\frac{g^2}{6}(-1,1,1),\qquad
 C=\frac16(-g,0,-g^3),\qquad
 \nb=\frac{(-1,1,1)}{\sqrt3},\qquad
 \nc=\frac{(-g^{-1},0,-g)}{\sqrt3}.
\end{equation}
The dot product is $\nb\cdot\nc=-1/3$, so $\cos\alpha=1/3$.
For a direct calculation of $A$, let
$\mathbf v=\nb\times(-\nc)$ and let $[\mathbf v]_\times$ be its
cross-product matrix. The shortest alignment is
\begin{equation}\label{eq:shortmatrix}
 A=I+[\mathbf v]_\times+
 \frac{[\mathbf v]_\times^2}{1+\cos\alpha}.
\end{equation}
Select $P_B=(-g/2,0,1/2)$ and $P_C=(-g/2,0,-1/2)$ in
\eqref{eq:outvectors}. Each face has circumradius $1/\sqrt3$, so
$\|\mathbf a\|^2=\|\mathbf b\|^2=1/3$. Direct substitution gives
\begin{equation}\label{eq:betacos}
 \boxed{\cos\beta=\frac{3w-1}{8},\qquad
 \sin\beta=\frac{\sqrt3(1+w)}8,\qquad
 \beta\simeq44.4775121859\degc.}
\end{equation}
The two displayed numbers have squared sum one and are positive, fixing
the quadrant. The three matching twists for face $8$ are
\begin{equation}\label{eq:correctedtwists}
 \boxed{\beta-120\degc,\qquad\beta,\qquad\beta+120\degc.}
\end{equation}
Their sine and cosine values are especially simple:
\begin{center}
\renewcommand{\arraystretch}{1.4}
\begin{tabular}{@{}lcc@{}}
\toprule
Twist & Cosine & Sine\\
\midrule
$\beta-120\degc$ & $1/4$ & $-\sqrt{15}/4$\\
$\beta$ & $(3w-1)/8$ & $\sqrt3(1+w)/8$\\
$\beta+120\degc$ & $-(3w+1)/8$ & $\sqrt3(w-1)/8$\\
\bottomrule
\end{tabular}
\end{center}
For the three faces with base $-\beta$, the matching set is the negative
of this set. These supply the opposite handedness of the same three
joint-orbit parameter families. The other class-$T$ face bases follow either
by the face symmetries or by the same substitution in~\eqref{eq:betarule}.

For the face-$8$ catalog representatives, every original twist is greater
than its matching replacement by
$45\degc-\beta\simeq0.5224878141\degc$. The nearest corresponding vertices
then differ by the positive distance
\begin{equation}\label{eq:mismatch}
 \frac{2}{\sqrt3}\sin\!\left(\frac{45\degc-\beta}{2}\right)
 \simeq0.00526492
\end{equation}
at edge length one. This is a finite geometric mismatch, not floating-point
roundoff. The implementation's use of the same positive $45\degc$ base
for all six class-$T$ faces also misses the negative-base family.
Nevertheless, all three literal catalog configurations remain well-defined
normal-aligned joint-centroid helices. Their parameter formulas in
Section~\ref{sec:icotables} are valid independently of exact vertex matching.

\subsection{Exact corrected parameters}
Retaining the names as labels for their corrected joining rules gives the
following replacements. The other 25 catalog entries are unchanged by this
vertex-matching correction. Edge length is one.
\begin{center}
\small\renewcommand{\arraystretch}{1.55}
\begin{tabular}{@{}lccc@{}}
\toprule
Name & $\tau$ & $\cos\theta$ & $\cos\phi$\\
\midrule
Quasi-planar & $\beta+120\degc$ & $\displaystyle \frac{- 3 w - 5}{12}$ & $\displaystyle \frac{3 g}{\sqrt{3 w + 17}}$\\
Icosacorkscrew & $\beta-120\degc$ & $\displaystyle - \frac{1}{6}$ & $\displaystyle - \frac{3}{\sqrt{14}}$\\
Big Icosacorkscrew & $\beta$ & $\displaystyle \frac{3 w - 5}{12}$ & $\displaystyle \frac{3}{g \sqrt{17 - 3 w}}$\\
\bottomrule
\end{tabular}
\end{center}
\begin{center}
\small\renewcommand{\arraystretch}{1.6}
\begin{tabular}{@{}lcc@{}}
\toprule
Name & $r$ & $d$\\
\midrule
Quasi-planar & $\displaystyle \frac{2}{\sqrt{3} \left(3 w + 17\right)}$ & $\displaystyle g^{3} \sqrt{\frac{2}{3 w + 17}}$\\
Icosacorkscrew & $\displaystyle \frac{\sqrt{15} g^{2}}{21}$ & $\displaystyle \frac{- g^{2}}{\sqrt{7}}$\\
Big Icosacorkscrew & $\displaystyle \frac{2 g^{4}}{\sqrt{3} \left(17 - 3 w\right)}$ & $\displaystyle g \sqrt{\frac{2}{17 - 3 w}}$\\
\bottomrule
\end{tabular}
\end{center}
\begin{center}
\small\renewcommand{\arraystretch}{1.3}
\begin{tabular}{@{}lrrrrr@{}}
\toprule
Name & $\tau$ (deg.) & $r$ & $d$ & $\theta$ (deg.) & $\phi$ (deg.)\\
\midrule
Quasi-planar & $164.477512$ & $0.048705$ & $1.230350$ & $167.3389$ & $4.4993$\\
Icosacorkscrew & $-75.522488$ & $0.482838$ & $-0.989524$ & $99.5941$ & $143.3008$\\
Big Icosacorkscrew & $44.477512$ & $0.769004$ & $0.713275$ & $81.8161$ & $54.6937$\\
\bottomrule
\end{tabular}
\end{center}

The source analog counts for these three rows are six each: three faces
with base $\beta$ and three with base $-\beta$, paired by handedness.
There are
\[
 3\cdot3+5\cdot4+7\cdot3+11\cdot5+19\cdot3=162
\]
distinct-face matching constructions across the five solids. Their counts
by representative, retaining the earlier names, are
\begin{center}
\small\renewcommand{\arraystretch}{1.2}
\begin{tabular}{@{}llr@{}}
\toprule
Solid & Representative multiplicities in catalog order & Total\\
\midrule
Tet & $6,3$ & $9$\\
Cub & $4,4,8,4$ & $20$\\
Oct & $3,6,6,3,3$ & $21$\\
Dod & $5,10,10,10,5,10,5$ & $55$\\
Ico & $3,6,6,6,12,3,6,6,6,3$ & $57$\\
\bottomrule
\end{tabular}
\end{center}
These multiplicities group the fixed-solid, fixed-edge joint-centroid orbits,
with handedness paired and with the original catalog's grouping of straight
orbits retained. Different twists about an opposite face can give the same
straight joint orbit while rotating the solid by different angles. The
multiplicities do not identify different solids or different ordered circular
orbits just because their containing circles happen to agree.

\section{Verification, scope, and further work}\label{sec:verification}
The supplemental exact catalog records both the literal 28-row input set
and the three matching corrections. A symbolic generator checks every
displayed $r^2,d^2,\cos^2\phi$ against the theorem, with the signs checked
separately, and obtains $\cos\theta$ directly. These checks reduce
identities in $\mathbb Q(\sqrt2,\sqrt3,\sqrt5)$ after clearing denominators;
they do not rely on agreement of decimal expansions. The exact vertex calculation
of $\cos\beta,\sin\beta$ is checked symbolically as well.

An independent standard-library Python program constructs the unit-edge
coordinate solids and their numbered faces, obtains rotations from
Rodrigues matrices, obtains axes from their antisymmetric parts in
nondegenerate cases, and measures radius using projected three-point
circumcircles. It compares the exact expressions with those constructions,
converts the lengths back to the original constructor scale, and checks all
four printed catalog columns. The five zero-angle markers are converted to
the stated straight-point convention for that comparison.

In the accompanying run, all 28 rows agree with the original table at its
three-decimal precision. The largest discrepancies against the rounded
entries are less than $0.000497$ in radius, $0.000453$ in advance,
$0.000476\degc$ in rotation, and $0.000444\degc$ in segment-axis angle.
The exact expressions and the independent numerical constructions agree
within $5\times10^{-16}$ for their algebraic outputs.

The vertex audit checks every corner under every placement. All 162
corrected distinct-face rules match to numerical precision; 18 rules from
the original implementation fail that check, precisely the three choices
on each of the six class-$T$ icosahedral faces. The maximum corrected
unit-edge vertex discrepancy is less than $6\times10^{-16}$ in the
accompanying run. A further 3,948 checks cover all ordered nonparallel face
pairs of the five solids at seven continuous twist values. Their largest
parameter discrepancy is less than $2\times10^{-15}$.

Numerical checks support the derivations and catch transcription errors;
they do not replace the geometric and algebraic proofs. Vertex matching
also does not establish the absence of collisions between nonadjacent
members of an infinite chain. The present results concern the prescribed
rigid placement and the orbit of its joint centers.

The equal-radius centroid hypothesis is the useful boundary of the present
position formulas. It permits generalization beyond the Platonic solids
whenever two selected joint centers lie at equal distances on their face
normals from a common point. Extending $r,d,\phi$ to arbitrary joint
positions requires the translational part of the placement in addition to
the two normals. The rotational half-angle law already survives that
extension. Further work can address collision-free realizations, exact
closure conditions, and geometrically appropriate illustrations of these
relationships.

\appendix
\section{Numerical unit-edge crosswalk}\label{app:numeric}
This table evaluates the original input twists, including the three literal
icosahedral choices. It is a crosswalk for reviewing the radical tables;
the exact expressions, rather than these rounded numbers, define the results.
Here $\theta$ is the full subunit rotation and straight rows use $d>0,\phi=0$.
\small
\renewcommand{\arraystretch}{1.1}
\begin{longtable}{@{}lrrrrr@{}}
\toprule
Name & $\tau$ (deg.) & $r$ & $d$ & $\theta$ (deg.) & $\phi$ (deg.)\\
\midrule
\endfirsthead
\toprule
Name & $\tau$ (deg.) & $r$ & $d$ & $\theta$ (deg.) & $\phi$ (deg.)\\
\midrule
\endhead
\midrule
\multicolumn{6}{r}{\emph{Continued on the next page}}\\
\endfoot
\bottomrule
\endlastfoot
Tetrahelix & $-120$ & $0.057735$ & $-0.316228$ & $131.8103$ & $161.5651$\\
Tetratorus & $0$ & $0.288675$ & $0.000000$ & $70.5288$ & $90.0000$\\
\addlinespace[0.4em]
Boxbeam & $-180$ & $0.000000$ & $1.000000$ & $180.0000$ & $0.0000$\\
Staircase & $-180$ & $0.000000$ & $0.707107$ & $180.0000$ & $0.0000$\\
Blockhelix & $-90$ & $0.235702$ & $-0.577350$ & $120.0000$ & $144.7356$\\
Cubatorus & $0$ & $0.500000$ & $0.000000$ & $90.0000$ & $90.0000$\\
\addlinespace[0.4em]
Octabeam & $-60$ & $0.000000$ & $0.816497$ & $60.0000$ & $0.0000$\\
Octaspiky & $-120$ & $0.104973$ & $-0.426401$ & $146.4427$ & $154.7606$\\
Octamedium & $-120$ & $0.115470$ & $-0.632456$ & $131.8103$ & $161.5651$\\
Octagear & $0$ & $0.288675$ & $0.000000$ & $109.4712$ & $90.0000$\\
Treestar & $0$ & $0.577350$ & $0.000000$ & $70.5288$ & $90.0000$\\
\addlinespace[0.4em]
Dodecabeam & $-108$ & $0.000000$ & $2.227033$ & $108.0000$ & $0.0000$\\
Dodecadoubler & $-144$ & $0.158056$ & $-1.128508$ & $161.3006$ & $164.5496$\\
The Alternater & $-144$ & $0.165307$ & $-1.867378$ & $149.5197$ & $170.3063$\\
Dodecashaft & $-72$ & $0.491851$ & $-0.760398$ & $129.6575$ & $130.5007$\\
Dodecagear & $0$ & $0.688191$ & $0.000000$ & $116.5651$ & $90.0000$\\
Dodecacorkscrew & $-72$ & $0.765348$ & $-1.534764$ & $93.0258$ & $144.1102$\\
Dodecadonut & $0$ & $1.801707$ & $0.000000$ & $63.4349$ & $90.0000$\\
\addlinespace[0.4em]
Pearlshaft & $-60$ & $0.000000$ & $1.511523$ & $60.0000$ & $0.0000$\\
Quasi-planar & $165$ & $0.047037$ & $1.230603$ & $167.7643$ & $4.3467$\\
Two Strands & $120$ & $0.130101$ & $0.474702$ & $159.4456$ & $28.3400$\\
Slow Twist & $120$ & $0.161109$ & $1.382980$ & $124.3094$ & $11.6407$\\
Rock Candy & $120$ & $0.194329$ & $0.789367$ & $146.4427$ & $25.2394$\\
Icosa Tree Star & $0$ & $0.288675$ & $0.000000$ & $138.1897$ & $90.0000$\\
Icosacorkscrew & $-75$ & $0.486993$ & $-0.986177$ & $99.2526$ & $143.0416$\\
Planar point cluster & $0$ & $0.534404$ & $0.000000$ & $109.4712$ & $90.0000$\\
Big Icosacorkscrew & $45$ & $0.763752$ & $0.719429$ & $82.0638$ & $54.3429$\\
The Wheel & $0$ & $1.978609$ & $0.000000$ & $41.8103$ & $90.0000$\\
\end{longtable}
\normalsize

\Needspace{12\baselineskip}
\section{Reproducibility and draft provenance}\label{app:provenance}
The catalog inputs are from Table~1 of \cite{read2022}. The implementation
audit uses repository revision
\begin{center}\small\texttt{8be3839de2204b966c8aedb7c28f6a6f5c18076e}.\end{center}
Here \texttt{main.js} supplies the constructor scales, face numbering, and
twist maps, and \path{js/segment_helix_math.js} supplies the adjoining
construction \cite{repository}. The supplement reconstructs the numbered
faces on exact coordinate solids; the independent numerical checks use
those radicals evaluated in double precision.

The package contains the following material:
\begin{itemize}
 \item \texttt{PlatonicHelicesClosedForms.tex}: self-contained LaTeX source.
 \item \texttt{catalog\_data.py}: catalog transcription, exact-coordinate
       constructions evaluated numerically, and numbered faces.
 \item \texttt{ExactCatalog.json}: exact expressions, fixed constants,
       source entries, and corrected entries.
 \item \texttt{verify\_closed\_forms.py}: independent numerical checks,
       requiring only Python's standard library.
 \item \texttt{build\_exact\_catalog.py}: optional symbolic regeneration
       using SymPy; generated table fragments are also supplied.
 \item \texttt{verification-results.txt}: output of the accompanying run.
\end{itemize}
The manuscript source embeds its tables and has no external image or
bibliography dependencies. It can be compiled with \texttt{pdflatex} twice
or with \texttt{latexmk -pdf}. The symbolic supplement is optional for
compilation and for running the independent verifier.

This author-review draft was developed with ChatGPT assistance in derivation,
calculation, verification code, and writing. It has not been submitted or
peer reviewed. No priority claim is made for the quaternion or screw-geometric
ingredients.

{\small
\begin{thebibliography}{9}
\bibitem{read2022}
R. L. Read, Calculating the segmented helix formed by repetitions of identical
subunits thereby generating a zoo of Platonic helices,
\emph{Mathematics} \textbf{10} (2022), 2533.
\href{https://doi.org/10.3390/math10142533}{doi:10.3390/math10142533}.

\bibitem{lord2002}
E. A. Lord, Helical structures: the geometry of protein helices and nanotubes,
\emph{Structural Chemistry} \textbf{13} (2002), 305--314.
\href{https://doi.org/10.1023/A:1015863923728}{doi:10.1023/A:1015863923728}.

\bibitem{dauter2018}
Z. Dauter and M. Jaskolski, On the helical arrangements of protein molecules,
\emph{Protein Science} \textbf{27} (2018), 643--652.
\href{https://doi.org/10.1002/pro.3356}{doi:10.1002/pro.3356}.

\bibitem{sadler2019}
G. Sadler, F. Fang, R. Clawson, and K. Irwin,
Periodic modification of the Boerdijk--Coxeter helix (tetrahelix),
\emph{Mathematics} \textbf{7} (2019), 1001.
\href{https://doi.org/10.3390/math7101001}{doi:10.3390/math7101001}.

\bibitem{onaka2024}
S. Onaka, Periodic arrangements of tetrahedra having appearances similar to
that of the Boerdijk--Coxeter helix,
\emph{Scientific Reports} \textbf{14} (2024), 18260.
\href{https://doi.org/10.1038/s41598-024-69108-w}{doi:10.1038/s41598-024-69108-w}.

\bibitem{repository}
Public Invention, \emph{Segmented Helices}, source repository,
\href{https://github.com/PubInv/segmented-helixes}{\texttt{PubInv/segmented-helixes}}.
Audit revision: \texttt{8be3839de2204b966c8aedb7c28f6a6f5c18076e}.
\end{thebibliography}
}
\end{document}
